\documentclass{beamer}
\usetheme{default}
\usecolortheme{seagull}
\usefonttheme{structurebold}
\usenavigationsymbolstemplate{}
\definecolor{csred}{rgb}{0.91,0.42,0.39}
\definecolor{csgreen}{rgb}{0.5,0.79,0.62}
\definecolor{csyellow}{rgb}{0.95,0.73,0.44}
\definecolor{cslblue}{rgb}{0.52,0.75,0.96}
\definecolor{csdblue}{rgb}{0.29,0.47,0.61}

\usepackage[utf8]{inputenc}
\usepackage{tikz}
\begin{document}

\addtobeamertemplate{frametitle}{\vskip-0.7ex}

\begin{frame}
\frametitle{\small{Minicurso}}
\begin{center}
	\vskip-0.3cm
	\begin{tikzpicture}
		\node[fill=csdblue, text=white] {
        	\raisebox{0.6cm}[1.5cm]{%
            \begin{minipage}{9cm}
			\centering \Large{\textbf{Computação Científica com Python}}\\\small{\textbf{com aplicações à Oceanografia e à Biologia}}
           	\end{minipage}
            }
         };
    \end{tikzpicture}
\end{center}
\vskip-0.3cm
\small{%
\textbf{Datas:} de 13/09 a 06/10, terças e quintas, das 18:30 às 20:00\\
\textbf{Local:} Laboratório de Informática do CFH \textemdash\ Sala 327 (30 vagas)
}

\scriptsize{%
\begin{itemize}
  \setlength\itemsep{0em}
	\item[$\bullet$] Introdução ao Python: instalação, configuração e a utilização do \emph{console}. Operações aritméticas e utilização de funções matemáticas básicas.
	\item[$\bullet$] Manipulação de listas e strings; vetores e matrizes usando SciPy e Numpy
	\item[$\bullet$] Tratamento de dados em arquivos; leitura e escrita em arquivos; Utilização de bibliotecas para leitura de arquivos \texttt{.csv}, \texttt{.txt}, e \texttt{.xls}
	\item[$\bullet$] Utilização da biblioteca \emph{PANDAS} para tratamento estatístico de dados.
	\item[$\bullet$] Gráficos com a biblioteca \emph{matplotlib}
	\item[$\bullet$] Resolução de problemas aplicados, incluindo a utilização de outras bibliotecas.
\end{itemize}
}

\small{%
\textbf{Inscrições e informações:} \texttt{gomicrobes@gmail.com} \textbf{até dia 08/09}
}
\begin{columns}
\column{1cm}
\column{1cm}
\tiny{\textbf{Apoio:}}
\column{2.5cm}
\begin{center}\tiny{Departamento de Matemática\\CFM - UFSC}\end{center}
\column{1.6cm}
\includegraphics[width=1.5cm]{brasao_1678.png}
\column{1.6cm}
\includegraphics[width=1.5cm]{GOM1BW.png}
\column{1.6cm}
\includegraphics[width=1.5cm]{brasao_UFSC.png}
\column{2cm}
\end{columns}
\end{frame}
\end{document}
